\documentclass{wyqart}
\input{macro}

\newcommand{\lasso}{\mathrm{Lasso}}
\newcommand{\yq}[1]{{\color{red} #1}}


\title{How to construct confidence sets after sparse regression}

\begin{document}

\maketitle

\section{Known results for fixed-sample sparse regression}

For observations $x_i\in\Rf^p$ and responses $y_i\in\Rf$,
let
$$
	X_n=\left(x_1^\top,\ldots,x_n^\top\right)^\top\in\Rf^{n\times p}
	\quad\text{and}\quad
	Y_n=\left(y_1,\ldots,y_n\right)^\top\in\Rf^n.
$$
Assume the sparse linear model
$$
	Y_n=X_n\beta^\star+\varepsilon_n,
$$
where $\beta^\star\in\Rf^p$ is the true coefficient vector and
$\varepsilon_n\in\Rf^n$ is the noise vector.
Write $[p]=\{1,\ldots,p\}$, and define
$$
	S^\star=\{j\in[p]:\beta_j^\star\neq 0\}
$$
as the true support.

A sparse-regression method such as the Lasso estimates $\beta^\star$ by
$$
	\widehat{\beta}_{\lasso,n}
	=
	\argmin_{\beta\in\Rf^p}
	\left\{
		\frac{1}{2}\norm{Y_n-X_n\beta}_2^2
		+
		\lambda\norm{\beta}_1
	\right\}.
$$
The tuning parameter $\lambda>0$ is often chosen by cross-validation.
An alternative constrained formulation, with radius $r>0$, is
\begin{align*}
	&\min_{\beta\in\Rf^p}
	&&\frac{1}{2}\norm{Y_n-X_n\beta}_2^2 \\
	&\text{subject to}
	&&\norm{\beta}_1\leq r.
\end{align*}

Sparsity can instead be imposed explicitly through an $\ell_0$ constraint:
\begin{align*}
	&\min_{\beta\in\Rf^p}
	&&\frac{1}{2}\norm{Y_n-X_n\beta}_2^2 \\
	&\text{subject to}
	&&\norm{\beta}_0\leq k.
\end{align*}
Here $k\in\{0,\ldots,p\}$ is an upper bound on the number of nonzero
coefficients.

\yq{Check \emph{Sparse High-Dimensional Regression:
Exact Scalable Algorithms and Phase Transitions} for a fast implementation
of the $\ell_0$-constrained estimator.}

For a generic sparse-regression procedure,
let $\widehat{S}_n\subseteq[p]$ denote the retained features.
The features in $\widehat{S}_n^c=[p]\setminus\widehat{S}_n$ are called the
kicked-out or excluded features.

\yq{Does the intended sparse-regression objective also include an
$\ell_2$ penalty, as in the elastic net?}

\yq{Identify and summarize the known fixed-sample guarantees for support
recovery or support containment.
Are they mostly asymptotic as $n\to\infty$ and conditional on assumptions
about the oracle parameter?}

\section{Anytime-valid sequential model comparison for linear regression}

For $S\subseteq[p]$, let $\Lf_S$ denote the class of linear models whose
coefficient vectors have support exactly $S$.
For example, consider the two models
$$
	y=\beta_1x_1+\beta_2x_2+\varepsilon
	\qquad\text{and}\qquad
	y=\beta_1x_1+\beta_3x_3+\varepsilon.
$$
They are represented by $\Lf_{\{1,2\}}$ and $\Lf_{\{1,3\}}$,
respectively.

A feature-level testing problem can be written as
\begin{equation}
	\Hc_0:\beta_3=0,
	\qquad
	\Hc_1:\beta_3\neq 0.
	\label{eq:test-one-coef}
\end{equation}
When $p$ is large, this is a comparison between the composite model classes
$$
	\{\Lf_S:3\notin S\}
	\qquad\text{and}\qquad
	\{\Lf_S:3\in S\}.
$$
The other included features and their coefficients are not fixed,
so both sides contain many possible linear models.

\yq{Ask S. Ariaus for}

\yq{Identify the precise existing anytime-valid result for OLS model
comparison.}

\yq{A possible first step is to construct an e-process for a suitably
specified simple-vs-simple comparison and then extend it to these composite
model classes.}

\section{Why pairwise model e-processes may be unsuitable for variable selection}

Even an optimal e-process for the comparison in
Equation~\eqref{eq:test-one-coef} would not by itself solve the variable-selection
problem.
With a large number of features, one would need to handle a large collection
of feature-level hypotheses or candidate models.

The comparisons also overlap.
For instance, comparing $\Lf_{\{1,2\}}$ with $\Lf_{\{1,3\}}$ concerns both
whether feature $2$ should be excluded and whether feature $3$ should be
included.
More generally, a single model contains many feature-level decisions,
and a single feature occurs in many candidate models.
It is therefore unclear how guarantees for separate model comparisons should
be combined into a guarantee for the feature set returned by a
sparse-regression algorithm.

\yq{Check whether the SMCS paper by Sebastian Arnold provides a useful
formulation of these overlapping comparisons.}

\section{Desired inferential guarantee}

In the conventional fixed-sample setting, $n$ is predetermined.
In the sequential setting, observations $(x_t,y_t)$ arrive online for
$t=1,2,\ldots$.
Let $X_t$ and $Y_t$ contain the first $t$ observations,
and let $\widehat{S}_t$ be the feature set retained by the sparse-regression
algorithm at time $t$.
A data-dependent stopping time is denoted by $\tau$.

The goal is to certify the features kicked out by time $\tau$.
More precisely, the desired guarantee is
$$
	\mathbb{P}\left(S^\star\subseteq\widehat{S}_\tau\right)
	\geq 1-\alpha.
$$
Here $\alpha\in(0,1)$ is the chosen error level,
and the probability covers the randomness in the sequential data and in the
procedure.
Equivalently, with probability at least $1-\alpha$,
every feature in $\widehat{S}_\tau^c$ is truly inactive.
This guarantee does not assert that every retained feature is active;
$\widehat{S}_\tau$ may still contain inactive features.

\section{Open questions and literature notes}

\begin{enumerate}
	\item Run the canonical sparse regression with $\ell_0$ control with fixed sample size $n$, how is the confidence sets constructed
	\item Try to make a example where the confidence interval is wrong, i.e., a true feature has been wrongly kicked out.
\end{enumerate}

\yq{How should an e-process be constructed for the feature set or exclusion
set returned by a sparse-regression algorithm?}

\yq{If the sparsity level $k=|S^\star|$ is known in advance,
can an explicitly constrained sparse-regression procedure be used to obtain
the desired guarantee for the excluded features?}

\yq{If the Lasso or elastic net can be run efficiently and already has
large-sample guarantees, how would an e-process
provide for the retained or excluded feature set?}

\yq{Literature to verify and summarize in a table or list:
work by Michael Lindon and Sebastian Arnold on model confidence sets,
and the Bayes-factor approach described by Ly et al.
\url{https://doi.org/10.1037/met0000454}.}

\end{document}
